% Perfect Square Condition for the Middle Area

\paragraph{Setup}
Let \(W_1\) and \(W_2\) be positive integers (not necessarily coprime).  
For positive integers \(n, m\) define
\[
M = nW_1 + mW_2 ,\qquad k = n + m .
\]
We form a square of side length \(M\); its total area is \(M^2\).  

The tiling algorithm fills in the boundary with tiles using \(4(k-1)\) rectangular tiles of size \(W_1 \times W_2\).  
The total area covered by these tiles is 
\[
A_{\text{tiles}} = 4(k-1) W_1 W_2 .
\]
The remaining uncovered region (the ``middle area'') is
\[
M = M^2 - A_{\text{tiles}}
    = (nW_1 + mW_2)^2 - 4(n+m-1)W_1W_2 .
\]

\paragraph{Special cases}
\noindent\textbf{Case 1: \(n = 1\).}  
Then \(M = W_1 + mW_2\) and \(k = 1 + m\).  Substituting into the expression for \(M\):
\begin{align*}
M &= (W_1 + mW_2)^2 - 4m W_1W_2 \\
  &= W_1^2 + m^2W_2^2 + 2mW_1W_2 - 4mW_1W_2 \\
  &= W_1^2 + m^2W_2^2 - 2mW_1W_2 \\
  &= (W_1 - mW_2)^2 .
\end{align*}
Hence \(M\) is a perfect square for any \(m \ge 1\).

\noindent\textbf{Case 2: \(m = 1\).}  
Then \(M = nW_1 + W_2\) and \(k = n + 1\).  Substituting:
\begin{align*}
M &= (nW_1 + W_2)^2 - 4n W_1W_2 \\
  &= n^2W_1^2 + W_2^2 + 2nW_1W_2 - 4nW_1W_2 \\
  &= n^2W_1^2 + W_2^2 - 2nW_1W_2 \\
  &= (nW_1 - W_2)^2 .
\end{align*}
Again \(M\) is a perfect square for any \(n \ge 1\).

\paragraph{Conclusion}
If either \(n = 1\) or \(m = 1\), the middle area \(M\) reduces to \((W_1 - mW_2)^2\) or \((nW_1 - W_2)^2\) respectively, which is always a perfect square.  
This holds for arbitrary positive integers \(W_1, W_2, n, m\).